\section{A norm anchor and a three-contact scalar certificate}

These changes start from the previously delivered 450-coordinate system.
They retain its inherited group, quaternion, height, pole-realization,
local-comparison, Lucas, and prime-pattern results. They change local
predicates and the input map, so their correctness is a transfer of the
soundness and completeness argument, rather than a birational equivalence
of the entire old and new systems.

\subsection{Use the existing quadratic field norm in the anchor}

Put $K=\mathbb Q(\rho)$, $\rho^2=2$, and
$N(a+b\rho)=a^2-2b^2$. Retain
$C_{\rm a}=3^{2^{25}}$, $S_0=C_{\rm a}^2$, and $B_0=4096C_0$.
Here $C_0$ is the rational coefficient of the product of the four root
abscissas. If $f_0+f_1\rho$ is the scalar reciprocal abscissa, replace
the physical anchor by
\[
 \Theta=N(f_0+f_1\rho)+3S_0N(B_0+X\rho)
 =f_0^2-2f_1^2+3S_0(B_0^2-2X^2).
\]
Keep the same conic
\[
 u^2+6v^2-2w^2-3\Theta^2+2S_0^2=0.
\]
The inherited quaternion criterion tests integrality of $\Theta/S_0$
at $2$ and $3$.

Since $2$ is a nonsquare in $\mathbb F_3$, the norm form is anisotropic
modulo $3$, and
\[
 v_3(N(a+b\rho))=2\min(v_3(a),v_3(b)).
\]
The summands of
\[
 \Theta/S_0
 =N(f_0/C_{\rm a}+(f_1/C_{\rm a})\rho)
       +3N(B_0+X\rho)
\]
have opposite valuation parity when nonzero. Integrality therefore
forces $f_0/C_{\rm a},f_1/C_{\rm a},B_0,X$ to be $3$-integral.
As in the preceding packed-anchor proof, this gives the old depth
condition and all four odd root indices. It then makes $B_0$ $2$-integral;
the scalar reciprocal coefficients divided by the odd integer $C_{\rm a}$
are already $2$-integral on the point domain.
All terms except $-6X^2$ are now $2$-integral. If $v_2(X)<0$, then
$v_2(-6X^2)=1+2v_2(X)<0$, which cannot cancel against those terms.
Thus $X$ is $2$-integral. Conversely these depth, parity and local-input
conditions make both displayed norm summands integral at both primes,
and the same quaternion criterion supplies conic witnesses.

The new anchor is consequently equivalent, on the retained point domain,
to the old depth condition, oddness of all four roots, and
$X\in\mathbb Z_{(2)}\cap\mathbb Z_{(3)}$. It removes the mixed term
$f_0f_1$ and expresses the entire physical quantity using the one existing
quadratic field norm. No new theorem about quaternion norm forms is needed.

\subsection{Make one finite contact a square fibre}

Let
\[
 A_0=2754=2\cdot3^4\cdot17,\qquad
 \beta=2401+A_0(2X+1)^2,
\]
and use
\[
 r=\beta(\beta-1)(32\beta+343),\qquad c=42.
\]
For every rational $X$, $\beta\ge2401$, hence $r>0$.
The exceptional equality $\beta=2401$ occurs exactly at $X=-1/2$.
It is excluded by the anchor's $2$-integrality conclusion before applying
any height theorem whose domain omits the contact set. Thus every solution
has $\beta\notin\{0,1,2401,-343/32\}$.

At every consumed prime $p>76176$, $A_0$ and $2$ are units. Therefore
\[
 v_p(\beta-2401)=2v_p(2X+1)
\]
is even, unless the expression is zero, which was just excluded globally.
If $v_p(\beta)<0$, the quadratic term in the expansion of $\beta$ dominates
its linear and constant terms, so $v_p(\beta)=2v_p(X)$ is even. Neither the
contact at infinity nor that at $2401$ contributes to $M_5(\beta)$.

The remaining three contacts $0,1,-343/32$ are pairwise distinct modulo
every consumed prime: their nonzero differences have numerator and
denominator prime factors only in $\{2,3,5,7\}$. If $\beta$ is integral
at such a prime, at most one of the three factors of $r$ has positive
valuation. Hence
\[
 p\mid M_5(\beta)\quad\Longleftrightarrow\quad v_p(r)\text{ is odd};
\]
in this case $v_p(r)>0$. If $\beta$ has negative valuation, all three
factors have that same even valuation, so $v_p(r)$ is even as well.
This checks both directions, although the scalar proof only needs the
forward implication.

The source scalar norm proof now applies verbatim to its valuation
recursion: $c=42$ is a unit at every consumed prime, and each odd contact
forces $v_p(a-1)\ge32$. Thus $M_5(\beta)^{32}\mid a-1$ and the inherited
height theorem gives $h(\beta)\le H_{19}|a-1|^{19/32}$.

\subsection{Fixed local correction and completeness}

For integral $X$, the odd integer $2X+1$ has square congruent to $1$ or $9$
modulo $16$. Since $A_0\equiv2\pmod{16}$,
\[
 \beta\equiv3\pmod{16},\qquad
 \beta\equiv2401\pmod{81},\qquad
 \beta\equiv4\pmod{17}.
\]
These congruences imply
\[
 v_2(r)=1,\quad r/2\equiv5\pmod8;
 \qquad
 v_3(r)=3,\quad r/27\equiv2\pmod3;
 \qquad r\equiv8\pmod{17}.
\]
For the $3$-adic calculation, the valuations of
$\beta,\beta-1,32\beta+343$ are respectively $0,1,2$, and their normalized
residues are $1,2,1$ modulo $3$.
Since
\[
 v_2(42)=v_3(42)=1,\quad42/2\equiv5\pmod8,
 \quad42/3\equiv2\pmod3,\quad42\equiv8\pmod{17},
\]
we obtain
\[
 42r>0,\qquad
 42r\in\mathbb Q_2^{\times2}\cap\mathbb Q_3^{\times2}
                 \cap\mathbb Q_{17}^{\times2}.
\]
The correction has a natural explanation: for $R(T)=T(T-1)(32T+343)$,
\[
 R(2401)=42\cdot102900^2.
\]
In fact a universal correction $1$ is impossible for this square-fibre
family: integers $X$ can make $2X+1$ arbitrarily $3$-divisible, making
$R(\beta)$ arbitrarily close to $R(2401)$, whose $3$-adic valuation is odd.

For integral $X$, $\beta\ge5155$ is positive and odd. In the unchanged
scalar norm tree, choose the shared point deeply at the support of $42r$,
then choose $[a-1]P$ deeper with the prescribed divisibility and anchor.
The same all-left constant leaf makes the root $W$ a unit at that support.
Thus $b=42rW^2$ has square reduced denominator and the three required local
squareclasses. The inherited pole-completeness lemma supplies witnesses.
Its continuity argument still gives a progression for the shared index.
No new quotient or point exception is introduced.

For soundness, once the inherited global argument gives $\beta\in\mathbb Z$,
the anchor has already made $X$ integral at $2$ and $3$. If an odd prime
$p\notin\{3,17\}$ divided the denominator of $X$, then
$v_p(\beta)=2v_p(X)<0$, a contradiction. At $17$ it would instead be
$v_{17}(\beta)=1+2v_{17}(X)<0$. Hence $X$ is integral everywhere.
The positive odd label and all subsequent scalar, product-height, Lucas,
prime-pattern and progression arguments are therefore valid with this
input map, subject to their unchanged inherited hypotheses.

\subsection{Keep the input gate adapted to the new label}

Materialize $Y=X^2+X$ as the one input-product coordinate, and use
\[
 \beta=11016Y+5155,\qquad X^2=Y-X.
\]
The defining row is $Y-X^2-X=0$. This is an invertible affine change of
the former $X^2$ coordinate over the fixed input $X$; it leaves the input
itself unchanged. Keep $Y$ protected. Then the anchor is implemented as
\[
 f_0^2-2f_1^2+50331648S_0C_0^2-6S_0(Y-X).
\]
All residuals remain quadratic even when $S_0$ is a constant-chain
coordinate. The input gate prevents the extra linear $X$ term from
propagating through every comparison circuit.

This arithmetic change does not by itself save a variable: the previous
four-factor scalar product already shared the quadratic leading part of
its two pair products. The combined norm-anchor and adapted-input variant
has the same $450$ variables and $419$ residuals as its baseline, with
$26921$ rather than $26922$ quartic monomials. These are measured outputs,
not a factor-count estimate.

\section{Independent review of homogeneous pole coordinates}

For a nonzero pole input $b$, put $B=17b$. The old equations are
\[
 Bp_{\rm old}^2=(B-e_{\rm old})(B+3e_{\rm old}),\qquad
 Bq_{\rm old}^2=(s_{\rm old}^2-H)(s_{\rm old}^2+3H),
\]
where $H=e_{\rm old}(17b-U^2)$. Set
\[
 e_{\rm old}=bE,\quad p_{\rm old}=bP,\quad
 s_{\rm old}=bS,\quad q_{\rm old}=b^2Q,\quad T=U^2/b.
\]
Then $H=b^2E(17-T)$ and the equivalent normalized equations are
\[
 17bP^2=(17-E)(17+3E),\qquad
 17bQ^2=(S^2-E(17-T))(S^2+3E(17-T)).
\]
The old residuals equal $b^2$ and $b^4$ times their respective normalized
residuals. The inverse coordinate maps divide only by $b$ and $b^2$.

For the three unprepared detector poles, $U=b$ and hence $T=b$. This
removes the formerly materialized $b^2$ coordinate in each pole module.
Here $b\ne0$ follows from the proved nonzero detector trace $w_0\ne0$
(an explicit guard in the preceding construction, supplied instead by
the positive-product lemma in the merged construction), together with
$2w_0n=w_0^2-2w_1^2$ and $be\delta=n^2$. Irrationality of
$\sqrt2$ gives $n\ne0$, whence $b\ne0$.
More precisely, in the merged construction the positive-$\lambda$ proof
gives $\sigma(W)\ge1$ at both real embeddings, hence $w_0\ge1$ and
$N(W)>0$, while $e\delta>0$. Therefore
$b=(N(W)/(2w_0))^2/(e\delta)>0$ before using any pole equation.

For the scalar pole, retain the actual old preparation variables $U,V$,
but replace its two squared coordinates by the single coordinate $T$:
\[
 T=b-V^2,\qquad U^2=bT.
\]
These are quadratic, and equivalent to $U^2+bV^2=b^2$ because $b\ne0$.
The old squared coordinates recover as $U^2=bT$ and $V^2=b-T$.
The expression $E(17-T)$ is quadratic, so the normalized pole uses only
three preparation coordinates $U,V,T$. No exceptional point is removed.
In this application $b=42rW^2>0$ even before using the anchor, because
$r>0$ for every rational input and the constant all-left norm leaf makes
$W>0$. Thus the nonzero hypothesis is unconditional on every arithmetic
solution. This proves the proposed saving of three detector coordinates
and one scalar preparation coordinate at the exact equation level;
the final compiler must still confirm the combined count.
